\begin{savequote}
\includegraphics[width=5cm]{./images/motal-aass/00800/neuro2.jpg}

Bild: Motal
\end{savequote}
\chapter
[\CO{[BAU]} Basisuntersuchungen]
{\CC{[BAU]}\\Basisuntersuchungen}
\label{chp:BAU}

\ChapterIntro
{%
\noindent Dieses Kapitel behandelt die Basisuntersuchungen.
}
{%
\begin{description}
  \item[Maintainer:] Sebastian Gabriel
  \item[Autoren:] Diverse
  \item[Reviewer:] --
  \item[Version:] {\DocVersion} ({\DocVersionInfo})
  \item[SHA1:] {\DocGitCommit}
\end{description}

{\TxTeilkapitelAASS}
}

\Layout{0}{\TxQuerverweise}
{~

\noindent\includegraphics[angle=270,width=\linewidth]{./images/leitsystem/Leitschema-PAM_BAU.pdf}
}{\begin{itemize}
  \item \EE{Beschreibung der Medizinprodukte}: \CO{[MP1, MP2]} 
  (\ref{chp:MP1}, \ref{chp:MP2}).
  \item \EE{Untersuchungen}:                 \CO{[BAU]} (\ref{chp:BAU}).
  \item \EE{Sanitätshilfemaßnahmen}:         \CO{[SAN]} (\ref{chp:SAN}).
  \item \EE{Anamnese und Patientengespräch}: \CO{[AUP]} (\ref{chp:AUP}).
\end{itemize}}{}{}{}{}



% \section{Blutdruck- und Herzfrequenzmessung}

\label{topic:rr-messung}

% Grundlagen
% 
% Fallstricke
% 
% Die Blutdruckmanschette soll nach Möglichkeit nicht auf der Seite
% 
% \begin{itemize}
%     \item eines Shuntarmes bei Dialysepatienten
%     \begin{itemize}
%         \item Ein Dialyse-Shunt ist ein künstlicher Kurzschluss zwischen einer
%         Armarterie und einer Armvene. Dadurch wird die Vene erweitert und sie
%         pulsiert, mann kann viel mehr Durchfluss für die Dialyse erreichen.
%         Die Nahtstelle zwischen Vene und Arterie ist jedoch sehr sensibel, wenn
%         gestaut wird kann sie aufplatzen 
%         \item ${\rightarrow}$ Gefahr einer massiven Shuntblutung
%     \end{itemize}
%     \item eines Armödems
%     \item einer Brust-OP mit Lymphknoten-Entfernung (\"Odemgefahr)
%     \item einer Halbseitenschwäche oder -lähmung
% \end{itemize}
% 
% angelegt werden. 
% 
% \Cave{Shunt-Arme sind tabu für die Blutdruckmessung!}
% 
% Bitten Sie im Praktikum am KTW Dialysepatienten höflich Ihnen ihren
% Shunt-Arm zu zeigen!
% 
% Maßnahmen
% 
% \begin{enumerate}
%     \item Evtl. Funktion des Blutdruckmessgerätes und des Stethoskops prüfen
%     \begin{itemize}
%         \item Kontrolle der Dichte des Ventils, der Verbindungsschläuche und
%         der Manschette (Ventil schließen, Manschette etwas aufpumpen und
%         Zusammenpressen, Manschette ganz entlüften)
%         \item Kontrolle der Gängigkeit des Ventils (es soll sich leicht auf
%         und zuschrauben lassen)
%         \item Kontrolle des Manometers auf korrekte Eichung: Die Nadel muss sich
%         im Ruhezustand genau bei 0\,mmHg befinden
%         \item Ohrstöpsel des Stethoskops zeigen schräg nach vorne
%         \item Trichterfunktion durch Berühren der Membran prüfen
%     \end{itemize}
%     \item Patientenlagerung
%     \begin{itemize}
%         \item Arm freimachen und entspannt auf einen Tisch legen in leichter
%         Flexion und Supination
%         \item Oberarm auf Herzhöhe
%     \end{itemize}
%     \item Auswahl der geeigneten Manschettengröße
%     \begin{itemize}
%         \item Manschette laut Herstellerangaben wählen z.B Standardmanschette:
%         Oberarmumfang max. 31 cm
%         \item Messwertverfälschung vermeiden! Der aufblasbare Teil der
%         Manschette soll 80\% des Oberarmumfangs umfassen, die Breite der
%         Manschette mindestens 1/3 des Oberarmumfangs betragen
%         \item Zu kleine Manschetten führen zum Messen zu hoher Werte, zu weite
%         Manschetten können nicht am Arm fixiert werden
%     \end{itemize}
%     \item Handhabung des Stethoskops
%     \item Bestimmen der Anlegestelle für Manschette und der
%     Auskultationsstelle durch Palpation
%     \begin{itemize}
%         \item Arteria brachialis medial in der Ellenbeuge aufsuchen
%     \end{itemize}
%     
%     \item Anlegen der Manschette
%     \begin{itemize}
%         \item Manschette muss auf unbekleideter Haut aufliegen
%         \item mäßig straff
%         \item Abstand zur Ellenbeuge 2,5 cm
%         \item aufblasbaren Teil über A. brachialis zentrieren
%         \item Schläuche nicht im Weg,
%         \item CAVE: Halbseitenschwäche/-lähmung, Shuntarm bei
%         Dialysepatienten, Armödem
%     \end{itemize}
%     \item Puls fühlen beim Aufpumpen
%     \begin{itemize}
%         \item Schnelles Aufpumpen
%         \item Tasten des Radialispulses mit Finger 2-4
%         \item Aufpumpen bis 30 mmHg über jenen Druck, ab dem kein Radialispuls
%         mehr getastet wird
%     \end{itemize}
%     \item Messvorgang
%     \begin{itemize}
%         \item Kopf des Stethoskops auf a. brachialis legen
%         \item Druck langsam ablassen 2 mmHg/sec
%         \item Systolischer Wert: Beginn der turbulenten Korotkoffschen
%         Strömungsgeräusche
%         \item diastolischer Wert: Verschwinden der Korotkoffschen Geräusche
%         \item Wenn Geräusche überhaupt nicht verschwinden (z.\,B. bei Kindern,
%         Schwangeren etc.), so gilt als diastolischer Wert der Messpunkt, an dem
%         die Geräusche deutlich leiser werden
%         \item Wenn Geräusche verschwunden sind, noch weitere 10mmHg langsam
%         ablassen -- wenn dann keine Geräusche mehr hörbar sind, kann die
%         Manschette zügig entleert werden
%         \item CAVE: auskultatorische Lücke
%     \end{itemize}
%     \item Protokollierung der Blutdruckwerte
% \end{enumerate}











\section{Sehen -- Hören -- Fühlen }
\label{topic:sehen-hoeren-fuehlen}
\Layout{0}{Unsere eigenen Sinne}{
Die Seh-, Hör- und Tastsinne sind bei der körperlichen
Untersuchung unverzichtbar. Zuerst wird der Patient und sein
Körper betrachtet. Wir versuchen, uns bereits bekannte Symptome wahrzunehmen.
Hierbei können uns viele augenscheinliche Sachen auffallen: zum Beispiel
\E{Fettleibigkeit}, \E{Zyanose}, \E{Blässe},
\E{Schweiß}, \E{Fremdkörper}, \E{Verletzungen}, \E{Fehlstellungen},
\ldots

Der nächste Schritt betrifft das Hören: Hier sind zum
Beispiel die Atemgeräusche einzuordnen, die man eventuell schon ohne
Hilfsmittel vernehmen kann (brodelndes Atemgeräusch beim
Lungenödem, pfeifendes Atemgeräusch beim Asthmaanfall,
röcheln bei Atemwegsverlegung, \ldots)

Wenn wir gesehen und gehört haben, können wir versuchen zu
fühlen: Wir können den (Radialis-) Puls fühlen, können
(stehende) Hautfalten beurteilen oder das Abdomen abtasten. Ebenso können wir
die ungefähre Temperatur fühlen. }{
\begin{itemize}
  \item Zyanose
  \item Blässe
  \item Schweiß
  \item Körperhaltung
  \item Fremdkörper
  \item Verletzungen, Fehlstellungen
  \item Atemgeräusche
  \item Puls
  \item Stehende Hautfalten
  \item Temperatur
  \item Durchblutung des Nagelbetts
  \item \ldots
\end{itemize}
}{}{}{}




\section{Vitalwerte}
\label{topic:vitalwerte}

\Layout{0}{\TxBeschreibung}{
Das Erheben von Vitalwerten ist eine Routinetätigkeit und
sollte bei jedem Notfallpatienten erfolgen. Es gibt nur wenig
Situationen bei denen man darauf verzichten kann (zum
Beispiel tobende Patienten). Zu den Vitalwerten zählt man

\begin{itemize}
\item die \EE{Herzfrequenz}, 
\item den \EE{Blutdruck} und 
\item die \EE{Atemfrequenz}.
\end{itemize}

Die Messung der Herzfrequenz erfolgt mittels Puls fühlen oder
mittels Pulsoxymeter. Beim Puls fühlen wird für
\EE{15~Sekunden} der Puls gefühlt und die Pulsschläge
gezählt. Anschließend wird das Ergebnis \EE{mit vier
multipliziert}, man erhält die Pulsschläge \EE{pro Minute}
(\nicefrac{x}{min}). Neben der Frequenz wird auch die
Regelmäßigkeit beurteilt, das heißt ob der Puls \E{rhythmisch} oder \E{arrhythmisch} ist.

Alternativ kann die Messung mittels des \E{Pulsoxymeters}
erfolgen (\prettyref{topic:pulsoxymetrie}). Dabei misst ein
Gerät die Herzfrequenz und zeigt diese an. Es ist jedoch zu
beachten, dass bei Notfallpatienten die Pulsoxymetrie oft
nicht funktioniert (z.\,B. Zentralisation bei Schock, kalte Finger,\ldots).

Die Messung des Blutdrucks ist unter \prettyref{topic:rr-messung} beschrieben.

Bei der Messung der Atemfrequenz wird die Anzahl der Atemzüge pro Minute gezählt.\index{Atemfrequenz}
}{
\begin{itemize}
  \item Herzfrequenz
  \item Blutdruck
  \item Atemfrequenz
\end{itemize}
}{}{}{}








\subsection{Messung der Herzfrequenz -- HF}
\label{topic:pulsmessung}
% TODO : Layout-Anpassung
\TODO{GABS}{2010-mm-dd}{Messung der Herzfrequenz:
Layout-Anpassung notwendig.} {Dieser Teil entspricht nicht dem Standard-Layout!}

 
\Definition{Als \T{Herzfrequenz} (\T{HF})
bezeichnet man die Anzahl der Schläge des Herzens pro Minute.} 

Sie wird mittels Fühlen des Pulses
ermittelt. \marginpar{\E{\scriptsize Die Blutgefäße werden im Detail unter
\prettyref{topic:blutgefaesse-wichtige} besprochen.}} 
Dieser wird idealerweise an der
\EE{Radialarterie} (A. \E{radialis}) bzw. bei nicht
ansprechbaren oder schockierten Patienten an der
\EE{Halsschlagader} (A. \E{carotis}) ertastet. Bei
Neugeborenen und Säuglingen kann man auch an der
\EE{Oberarmarterie} (A. \E{brachialis}) fühlen. Gemessen
wird immer mit zwei bis drei Fingern. Nie jedoch mit dem Daumen, da hier meistens nur der eigene Puls gefühlt wird und nicht der des
Patienten. Wenn an der Halsschlagader (A. \E{carotis}) die
Herzfrequenz ermittelt wird, darf nie auf beiden Seiten
gleichzeitig gemessen werden, da sonst die Gefahr besteht,
dass es zu einer Unterversorgung des Gehirn durch Blut kommen
kann.

\Layout{2}{Durchführung}{
Es wird an der Radialarterie (A. \E{radialis}) der Puls mit
zwei bis drei Fingern ertastet. Nun werden 15 Sekunden lang
die Schläge gezählt. Multipliziert man den Messwert mit vier,
erhält man  die Herzfrequenz (HF) des Patienten. Der
Normalwert liegt beim Erwachsenen zwischen 60 und 100
Schlägen pro Minute. Dokumentiert wird sie mit einem HF
(z.\,B.: HF\,=\,70)
}{
}{}{}{}


\subsection{Messung des Blutdrucks -- RR}
\label{topic:blutdruckmessung}

\Layout{0}{\TxBeschreibung}{
Der Blutdruck ist jener Druck, welcher in unseren Blutgefäßen
herrscht. Wenn wir vom \emph{"`Blutdruck"'} sprechen beziehen wir
uns auf den Blutdruck in den
Arterien\footnote{\emph{Arterie:} Vom Herzen wegführendes
Blutgefäß.}. Er ist in Kombination mit der Herzfrequenz
einer der wichtigsten Messwerte im
Rettungsdienst.

Im Idealfall werden zwei Werte ermittelt: Der
\EE{systolische Wert} ("`obere Wert"') ist jener Druck im
Blutgefäßsystem, der  während der \E{Auswurfphase} des Herzens, in der sich
die Herzkammern anspannen, zusammenziehen und Blut auswerfen, herrscht. Der
\EE{diastolische Wert} ("`untere Wert"') ist jener Druck im
Blutgefäßsystem, der während der \E{Entspannungsphase} des Herzens,
in der sich die Herzkammern mit
Blut füllen, herrscht. 

Die \EE{Einheit}, in der der Blutdruck angegeben wird, ist
\T{Millimeter Quecksilbersäule}
(\T{mm\,Hg}).\index{Millimeter Quecksilbersäule}\index{mm\,Hg}\footnote{\ce{Hg}
ist das chemische Symbol für Quecksilber.}
}{
\begin{itemize}
  \item Druck in den Arterien
  \begin{itemize}
    \item systolischer Wert (Auswurfphase des Herzens)
    \item diastolischer Wert (Entspannungsphase des Herzens)
  \end{itemize}
  \item Einheit \MmHg
\end{itemize}
}{}{}{}

\Layout{0}{Arten der Blutdruckmessung}{%
Der Blutdruck kann auf verschiedene Arten gemessen werden. Üblich ist die
unblutige, \E{indirekte Methode}, mittels einer an der Extremität angelegten
Druckmanschette\ZFootnote{Mittels einer Druckmanschette wird gemessen, bei
welchem Druck ein mit einem Stethoskop hörbarer, durch Pulstasten fühlbarer
oder durch sonstige technische Hilfsmittel messbarer Blutfluss stattfindet.
Dies erlaubt Rückschlüsse auf den Druck im Blutgefäß.}. Dazu gibt es unterschiedliche
Messgeräte, grundsätzlich unterscheidet man zwischen elektronischen und rein
mechanischen Messgeräten.

\EE{Elektronische Messgeräte}  werden oft von Patienten zur Selbstkontrolle des
Blutdrucks verwendet. Diese Geräte sind jedoch vergleichsweise fehleranfällig
und sind für einen professionellen Einsatz nicht geeignet. 

Für den
professionellen Einsatz gibt es spezielle
elektronische Blutdruckmessgeräte, welche zum Beispiel in
Überwachungsmonitoren zum Einsatz kommen. Allerdings sind auch diese Geräte 
mitunter fehleranfällig, wenn die Messbedingungen nicht optimal
(z.\,B. bei Erschütterungen) sind.

% Der Blutdruck kann mittels elektronischer
% Blutdruckmessgeräte ermittelt werden. Diese sind aber im
% Rettungsdienst nicht in Verwendung, da es immer wieder zu
% Gerätefehler kommen kann. Außerdem können die Batterien leer
% werden. Solche Geräte haben meistens die Patienten für die
% tägliche Überwachung ihres Blutdrucks. Leider bedienen aber
% viele die Geräte falsch und fürchten sich dann
% lebensbedrohlich krank zu sein.

In der Praxis wird die Blutdruckmessung mit \EE{mechanischen Messgeräten}
durchgeführt. Dabei wird eine aufblasbare Manschette mit angeschlossenem Dosenmanometer und
Pumpballen mit einem Ablassventil benötigt. Im Idealfall wird auch noch ein
Stethoskop verwendet.

\begin{center}
     
%   \includegraphics[width=4cm]{./images/noauth/GER-rr-elektrisch-oberarm.png}
%   
%    \includegraphics[width=4cm]{./images/noauth/GER-rr-elektrisch-unterarm.jpg}
%    
%    \includegraphics[width=4cm]{./images/noauth/GER-rr-messgeraet.jpg}
\TODO{GABS}{2010-04-27}{Blutdruckmesser: Bilder aus Serie
Pallinger}{}
\end{center}
}{
\begin{itemize}
  \item unblutige, indirekte Messung
  \begin{itemize}
    \item mechanische Messung
    \item elektronische Messung
  \end{itemize}
  \item invasive Messung
\end{itemize}
}
{}
{}
{}


\subsubsection{Messung nach Riva-Rocci}

\Definition{Die Blutdruckmessung nach
Riva-Rocci\ZFootnote{Benannt nach dem Erfinder Scipione
Riva-Rocci.} ist eine mechanische Messung mittels einer Druckmanschette.}

Umgangssprachlich wird die Blutdruckmessung nach
Riva-Rocci als \emph{"`RR-Messung"'}, bzw. der ermittelte Wert als \emph{"`RR"'} bezeichnet.

Bei der RR-Messung gibt es zwei Möglichkeiten zu
einem Wert zu kommen. Die genauere ist die
\EE{auskultatorische Methode} (durch hören von
Strömungsgeräuschen). Die zweite Möglichkeit ist die
\EE{palpatorische Methode} (durch tasten des Pulses). 

\Fazit{Der diastolische Wert kann nur
bei der auskultatorischen Methode ermittelt werden!} 

%%%%%%%%%%%%%%%%%%%%%%%%%%%%%%%%%%%%%%%%%%%%%%%%%%%%%%%%%%%
% TODO : Neuerstellung
\TODO{KOCH}{yyyy-mm-dd}{RR, Bild von Arm mit eingez. Arterie: Neuerstellung von
Foto notwendig}{}
%%%%%%%%%%%%%%%%%%%%%%%%%%%%%%%%%%%%%%%%%%%%%%%%%%%%%%%%%%%


\Layout{0}{\T{Auskultatorische Methode} -- Vorgehen}
{%
Die Manschette sollte faltenfrei und nicht über Kleidung, in der
Mitte des Oberarms
angelegt werden und sich, wenn möglich, auf
\E{Herzhöhe} befinden\ZFootnote{Wenn die Manschette nicht
auf Herzhöhe liegt, wird ein zu hoher oder zu
niedriger Druck gemessen.} Die Mitte der Luftkammer soll
über der Oberarmarterie zu liegen kommen. Manche Manschetten haben einen Markierungspfeil für die Oberarmarterie, der Pfeil soll dann oberhalb dieser zu liegen kommen. Man achte auf die Lage der Innen- bzw. Außenseite!

Nun wird
der Puls an der Radialarterie (A. radialis) getastet. Die Manschette wird
solange aufgepumpt, bis der Puls nicht mehr getastet werden kann. Dann werden
noch \VonBis{20}{30}\MmHg weiter aufgepumpt. Das Stethoskop wird in der Ellenbeuge auf
die Oberarmarterie (A. brachialis, Verlauf eher zur Körpermitte) angelegt. Die
Luft wird langsam aus der Manschette abgelassen. 

Der Wert, bei dem man
erstmalig ein Pochen \Z{(Korotkowsche Geräusche)} hört, gilt
als \EE{systolischer Wert}. Die Luft wird weiter abgelassen. Der Wert, bei dem das
Pochen verschwindet, ist der \EE{diastolische Wert}\footnote{Es kann
vorkommen, dass das Pochen nie verschwindet, sondern nur
plötzlich leiser wird. Dann gilt der Druck, bei dem das
Pochen \emph{deutlich} leiser wird, als diastolischer
Wert.}.

Die
Werte werden auf die nächste Fünferstelle abgerundet (z.\,B.:
\RRmmHg{122}{82} wird abgerundet auf \RRmmHg{120}{80}). Dokumentiert
wird der Blutdruck mit RR (z.\,B.: RR~\RRmmHg{125}{75}). Der
Normalwert beim Erwachsenen liegt ca. bei RR~\RRmmHg{120}{80}
(\prettyref{topic:blutdruck}).
}{
\begin{itemize}
  \item Manschette faltenfrei auf Herzhöhe anlegen
  \item Luftkissen  mittig über Oberarmarterie
  \item Markierungspfeil beachten
  \item Radialarterie tasten
  \item Manschette aufpumpen bis Radialispuls nicht mehr tastbar ist
  +~\VonBis{20}{30}\MmHg
  \item Stethoskop über Oberarmarterie anlegen
  \item Druck langsam ablassen
  \item Erstes Pochen: systolischer Wert
  \item Letztes Pochen: diastolischer Wert
  \item Werte auf nächste 5er-Stelle abrunden
\end{itemize}

}{}{}{}




\begin{figure}[H]
    \begin{center}
    \includegraphics[width=\linewidth]{./images/hirtler-aass/Riva-Rocci_edited.pdf}
   
    \Captionof{figure}{Prinzip
    der auskultatorischen RR-Messung: In Abhängigkeit vom außen anliegenden
    Manschettendruck kommt es zu einem unterschiedlich starken Blutfluss und dadurch entstehenden, hörbaren, Strömungsgeräuschen.}
    \end{center}
\end{figure}


\Layout{0}{\T{Palpatorische Methode} -- Durchführung}
{%
Die Manschette sollte faltenfrei und nicht über Kleidung, in der
Mitte des Oberarms
angelegt werden und sich, wenn möglich, auf
\E{Herzhöhe} befinden\ZFootnote{Wenn die Manschette nicht
auf Herzhöhe liegt wird ein zu hoher oder zu
niedriger Druck gemessen.} Nun wird der Puls an der
Radialarterie (A. radialis) ertastet. Die Manschette
wird solange aufgepumpt, bis der Puls nicht mehr getastet
werden kann, dann wird der Druck um weitere 20\MmHg
erhöht. Die Finger bleiben auf der Radialarterie (A.
radialis). Die Luft wird langsam aus der Manschette abgelassen. Der Druck, bei dem der
Puls wieder fühlbar ist, ist der \EE{systolische} Wert. 

\EE{Der diastolische Wert kann mit
dieser Methode nicht ermittelt werden.} In der
Notfallmedizin ist der systolische Wert jedoch ohnehin in der
Regel der wichtigere.

\bigskip

\bigskip

\bigskip

\bigskip

\bigskip
}{
\begin{itemize}
  \item Manschette faltenfrei auf Herzhöhe anlegen
  \item Luftkissen  mittig über Oberarmarterie
  \item Markierungspfeil beachten
  \item Radialarterie tasten
  \item Manschette aufpumpen bis Radialispuls nicht mehr tastbar ist
  +~\VonBis{20}{30}\MmHg
  \item Druck langsam ablassen, dabei Radialispuls weiter versuchen zu tasten
  \item Radialispuls wieder fühlbar: systolischer Wert
  \item Auf nächste 5er-Stelle abrunden
  \item \E{Kein} \E{diastolischer} Wert!
\end{itemize}
}{}{}{}



\Layout{60}{Was ist besonders zu beachten?}
{%
Die Blutdruckmanschette soll nach Möglichkeit \EE{\E{nicht} auf der Seite
\ldots}
\begin{inparaenum}
	\item eines \EE{Shuntarmes}\footnote{Ein (Dialyse-)\EE{Shunt} ist ein
	künstlicher Kurzschluss zwischen einer Armarterie und einer Armvene. Dadurch wird
		die Vene erweitert und sie pulsiert. Mann kann viel mehr
		Durchfluss für die Dialyse erreichen. Die Nahtstelle
		zwischen Vene und Arterie ist jedoch sehr sensibel, wenn
		gestaut wird kann sie aufplatzen $\rightarrow$ Gefahr
		einer massiven Shuntblutung.} bei Dialysepatienten
		\item eines Armödems, z.\,B. nach einer
		Brustamputation, oder
		\item einer Halbseitenschwäche oder -lähmung 
\end{inparaenum} 
angelegt werden.

\Cave{Shunt-Arme sind tabu für die Blutdruckmessung!}
}{
\begin{itemize}
  \item \EE{Nicht} auf Seite eines/r \dots{}
  \begin{enumerate}
    \item \EE{Shuntarmes}
    \item Armödems
    \item Halbseitenschwäche/-lähmung
  \end{enumerate}
\end{itemize}

}{./images/sign_cardive.jpg}{}{Gabriel}







\Z{
\Layout{2}{Weitere Bemerkungen}
{\index{NIBD|Z}\index{NIBP|Z}
Der "`Nicht-invasive Blutdruck"' wird auf manchen Überwachungsmonitoren mit
\T{NIBD} oder \T{NIBP} (\Lang{engl.} Non-Invasive Blood
Pressure) abgekürzt. }{}{}{}{}
}



\section{Apparative Untersuchungen}

\subsection{Pulsoxymetrie}
\label{topic:pulsoxymetrie}
% TODO: Pulsoxymetrie: Vergleich mit Bettwache einfügen [GABS] 
% TODO: Pulsoxymetrie: Sättigung sackst rasch ab, bezug zu Gas-Partialdrücken
% einfügen [GABS]

\Layout{2}{\TxBeschreibung}
{%
Die Pulsoxymetrie ist eine einfache apparative
Untersuchungsmethode, welche mithilfe eines
\E{Pulsoxymeters} durchgeführt wird. 
Sie dient zur nicht-invasiven Überwachung der
\T{Sauerstoffsättigung} (\T{Sp\ce{O2}}) des Hämoglobins\footnote{Das Hämoglobin ist der für die Sauerstoffbindung verantwortliche Stoff in den roten Blutkörperchen (\prettyref{topic:haemoglobin}).} im kapillären
Blut. Dies geschieht mittels Photorezeptoren und
Infrarotlicht, mit deren Hilfe das Verhältnis der mit
Sauerstoff beladenen roten Blutkörperchen zu den nicht
beladenen bestimmt wird. 

Als weiteren Wert zeigt das Pulsoxymeter die
\EE{Pulsfrequenz} an.
\cite{LoefflerBiochemie7,ArbeitstechnikenAZ1,Soehnke200607}

}{
% TODO Slide fehlt.
}{}{}{}

\Layout{2}{Durchführung}
{%
Es wird eine Klemme mit einer Lichtquelle
und einem Sensor auf den Finger des Patienten
geklippt. Das Licht wird vom
 Sensor aufgefangen, das Gerät kann aus den daraus gewonnenen Informationen auf den Sauerstoffgehalt des Blutes und die Pulsfrequenz schließen. 

Dies setzt jedoch
voraus, dass der Patient in der Peripherie über eine
ausreichende Durchblutung verfügt, da das Gerät sonst keine
Informationen gewinnen kann (mangels Blut zur
Beurteilung). 
}{}{}{}{}


\Fazit{Eine ausreichende Blutversorgung der Körperstelle,
an der der Sensor angebracht ist, ist Voraussetzung für die
Pulsoxymetrie.}

Das Pulsoxymeter eignet sich sehr gut für das Monitoring,
ersetzt aber niemals das wachsame Auge auf den Patienten!

\Layout{0}{Ermittelte Werte}
{%
Die Sauerstoffsättigung wird in Prozent
angegeben und wird mit \T{Sp\ce{O2}} abgekürzt. Der Normalwert
liegt zwischen \EE{\VonBis{96}{100}\PC*}. Bei Patienten mit
COPD findet man jedoch oft sehr niedrige Werte (z.\,B.
90\PC*), ohne dass die Patienten Beschwerden haben. }{
% TODO Slide fehlt.
% \begin{itemize}
%     \item \Z{Point}
% \end{itemize}
\begin{itemize}
  \item Normalwert: \VonBis{96}{100}\PC*
  \item COPD-Patient: auch Werte < 90\PC* möglich ohne
  Beschwerden
\end{itemize}
}{}{}{}


\Layout{2}{Beurteilung}
{
Bei der Interpretation des
$SpO_2$ muss man sehr vorsichtig sein! Auch ein Normalwert bedeutet nicht, dass der Patient
genug Sauerstoff hat, beziehungsweise, dass keine
Sauerstoffgabe erfolgen soll. Die Sauerstoffgabe richtet
sich immer nach dem klinischen
Zustandsbild des Patienten und dessen Diagnose, und nur in
Ausnahmefällen (z.\,B. Asthma bronchiale,
\prettyref{topic:asthma-bronchiale}) nach dem
Sp\ce{O2}-Wert. Der Sp\ce{O2}-Wert soll grundsätzlich im
Verlauf beurteilt werden.
\cite{Soehnke200607}

\Fazit{Behandle den Patienten und nicht das Gerät!}



Bei der Beurteilung des Wertes muss auch berücksichtigt
werden, \E{ob und wieviel \EE{Sauerstoff} der Patient
erhält}. Eine Sauerstoffsättigung von 95\PC* bei einem Patienten ohne
zusätzliche Sauerstoffgabe ist nicht beunruhigend,
wohingegen derselbe Wert bei einer \E{Sauerstoffgabe} von
15{\LPerMin*} ein Hinweis für eine schwere Störung der Atmung
darstellt. 

\Fazit{Die Beurteilung der Sauerstoffsättigung muss
berücksichtigen, ob und wieviel Sauerstoff der Patient
erhält.} 
}{
% TODO Slide fehlt.
}{}{}{}




\Layout{0}{\TxGefahren{} und Fehlerquellen}
{%
Eine besondere Fehlerquelle ist die
Kohlenmonoxid-(\ce{CO})-Vergiftung, bei ihr wird nicht nur
der Patient schweinchenrosa obwohl er eigentlich zu wenig
Sauerstoff hat, auch das Pulsoxymeter verwechselt
Kohlenmonoxid mit Sauerstoff und zeigt falsch hohe Werte an!

Voraussetzung für eine korrekte Messung ist eine ausreichende
Durchblutung der entsprechenden Messstelle, z.\,B. durch Kälte,
Schock oder auch eine Blutdruckmessung kann die Messung
gestört oder unmöglich gemacht werden. Weitere Hindernisse
sind z.\,B. kalte Finger (verringerte Durchblutung) oder
Nagellack (Licht kann den Nagellack nicht durchdringen,
bzw. dessen Farbe wird verfälscht.) Mit
Nagellack-Entferner kann der Nagel eventuell
gereinigt werden. Alternativ kann der Sensor auch
um 90\textdegree gedreht werden. 

Von der Seite einstrahlendes Licht kann ebenso das Ergebnis
verfälschen, da es sich ja um eine lichtabhängige
Messmethode handelt.
}{
\begin{itemize}
	\item \ce{CO}-Vergiftungen (falsch hohe Werte)
	\item Minderdurchblutung der Peripherie; Entweder durch Zentralisation beim Schock oder Unterkühlung.
	\item Nagellack
	\item Künstliche Fingernägel
	\item Lichteinstrahlung von der Seite
\end{itemize}
}{}{}{}




\subsection{Blutzuckermessung}
\label{topic:blutzucker-messung}

\Layout{0}{Interpretation und Normalwert}
{
Der Blutzuckerspiegel wird in \E{Milligramm pro Deziliter}
(\MgDl\Z{, entspricht \MetricUnit{mg\,\%}}) gemessen. Der
Wert ist abhängig von der letzten Mahlzeit des Patienten.
Nach der Nahrungsaufnahme steigt der Blutzuckerspiegel
stark an und fällt dann wieder auf einen "`nüchternen
Normalwert"' ab. Der normale
\EE{Nüchtern-Blutzuckerspiegel} bewegt sich zwischen
\EE{80\,--\,120\MgDl*}. Im Rettungsdienst
kann der Normalbereich jedoch großzügiger angesetzt werden,
da bei Werten zwischen 60 und 200\MgDl* keine
Notfallsymptome auftreten.

}{
\begin{itemize}
    \item Abhängig von letzter Mahlzeit!
    \item Normalwert: 80\,--\,120\MgDl*
    (nüchtern)
    \item Bereich ohne Notfallsymptome:
    60\,--\,200\MgDl*
%     \item Notfallrelevanter Hypoglykämie: <
%     60\MgDl*
%     \item Notfallrelevante Hyperglykämie: >
%     200\MgDl*
\end{itemize}
}{}{}{}

\Layout{2}{Material \& Personal}
{
\begin{enumerate}
  \item 1 Person
  \item Handschuhe
  \item Kontamed-Box
  \item Alkoholtupfer
  \item Trockene Tupfer
  \item Blutzuckermessgerät
  \item Blutzuckermessstreifen
  \item Lanzette
\end{enumerate}
}
{}
{./images/gabriel-sebastian-aass/bz-wert_1214-00800.jpg}
{"`90\MgDl*."'}
{GABS}


\Layout{2}{Durchführung}
{
%\begin{multicols}{2}
\begin{enumerate}
  \item Finger auswählen: möglichst von der nicht dominanten
  Hand.
  \item Einstichstelle bestimmen: Am Fingerendglied,
  seitlich (Rand!), nicht in die Fingerkuppe, nicht zu nahe
  am Gelenk. Außenseite des Ringfinger bevorzugen.
  \item Stelle desinfizieren: Mit Alkoholtupfer abwischen
  und gemäß Einwirkzeit einwirken lassen.
  \item Blutzuckermessstreifen halb in das Gerät schieben,
  aber noch nicht einrasten lassen.
  \item Mit einer schnellen, beherzten Bewegung den
  Patientenfinger mit der Lanzette stechen und sofort
  herausziehen \E{("`Rein-raus-Methode"')}.
  \item Lanzette in die \E{Kontamed-Box} entsorgen.
  \item Ersten Bluttropfen mit trockenem Tupfer abwischen.
  \item Messstreifen ganz in das Gerät schieben, das Gerät
  schaltet sich ein und führt einen Selbsttest durch. Wenn
  im Display ein Tropfen angezeigt wird ist das Gerät
  messbereit.
  \item Den Messstreifen mit dem Gerät in den Bluttropfen
  tauchen, das Blut steigt auf\footnote{Kapillarwirkung}.
  \item Beim Piepton ist genug Blut in die Testkammer
  eingeströmt und das Gerät beginnt mit der Messung und kann
  entfernt werden. 
  \item Mit einem trockenen Tupfer auf die Einstichstelle
  drücken (lassen).
  \item Ein neuerlicher Piepton zeigt an, dass die Messung
  beendet wurde und der Wert vom Display abgelesen
  werden kann. 
  \item Messstreifen entsorgen\footnote{Entsorgung gemäß
  Abfallentsorgungsplan. \opt{ORGASBFLO}{\ASBFLO{} Der
  Messstreifen kann in den normalen Müll entsorgt werden.}}.
\end{enumerate}
%\end{multicols}
}{

}{}{}{}

\begin{table}[H]
\begin{WidePar}
\Captionof{table}{Bilderserie: Blutzuckermessung}
\begin{tabularx}{\linewidth}{XXX}
\includegraphics[width=\linewidth]{./images/pallinger-christoph-aass/Blutzucker_32796_export_01312.jpg}
\Captionof{figure}{Material \SourcePicture{Ch. Pallinger}}
&
\includegraphics[width=\linewidth]{./images/motal-aass/00800/neuro6.jpg}
\Captionof{figure}{BZ-Messung \SourcePicture{Motal}}
&
\includegraphics[width=\linewidth]{./images/gabriel-sebastian-aass/bz-wert_1214-00800.jpg}
\Captionof{figure}{"`90\MgDl*."' \SourcePicture{GABS}}
\\\bottomrule
\end{tabularx}
\end{WidePar}
\end{table}



\subsection{Körpertemperatur}

Die Körpertemperatur kann an unterschiedlichen Stellen und
mit unterschiedlichen Mitteln bestimmt werden. 

\Layout{0}{Messarten}
{%
Mit dem
\EE{herkömmlichen Fieberthermometer}\footnote{Unter
herkömmlichen Fieberthermometern versteht man heute
elektronische Messgeräte. Quecksilberthermometer werden
wegen des darin enthaltenen, giftigen Schwermetalls
Quecksilber nicht mehr verwendet!} kann die Temperatur unter der \E{Achsel} (\E{axillär}) oder im \E{Enddarm} (\E{rektal}, vor allem bei Kleinkindern) bestimmt werden. Die Messung im Mund ist in der Praxis eher unüblich. Es ist bei der Messung eine Schutzhülle zu verwenden. 

Mit dem \EE{Ohrthermometer} kann mittels Infrarot die
Temperatur im \E{Mittelohr} bestimmt werden. Es sind die
entsprechenden \E{Schutzkappen} zu verwenden.

Die Temperatur im Enddarm und im Mittelohr bezeichnet man
als \T{Körperkerntemperatur}, diese beträgt um die
37\,\textcelsius. Die unter der Achsel gemessene Temperatur
kann ca. \nicefrac{1}{2}\,\textcelsius{} niedriger sein.
}{
\begin{itemize}
    \item nach dem Ort der Messung
    \begin{itemize}
      \item axillär
      \item rektal
    \end{itemize}
    \item nach dem Gerätetyp
    \begin{itemize}
      \item herkömmliches (elektronisches) Thermometer
      \item Ohrthermometer
    \end{itemize}
\end{itemize}
}{}{}{}





\subsection{EKG}
\label{topic:ekg}
\index{EKG}

Das \E{Elektrokardiogramm} (\Abk{EKG}) zeigt mittels eines
entsprechenden Gerätes die elektrische Herzaktivität des
Reizbildungs- und Reizleitungssystems an. Dazu werden auf
dem Körper nach einem bestimmten Muster Elektroden
aufgeklebt und mittels eines Kabels mit dem \Abk{EKG}-Gerät
verbunden. Im Rettungsdienst werden fast ausschließlich
Klebeelektroden für den einmaligen Gebrauch verwendet. 

\AuthNotes{GABS: Sollte im Defi-Teil durchgenommen
werden?:} \Z{EKG -- Ablauf einer Erregung

\begin{itemize}
    \item P-Welle $\approx$ Vorhoferregung
    \item PQ-Zeit $\approx$ AV-Überleitung
    \item QRS-Komplex $\approx$ Kammererregung
    \item ST-Strecke + T-Welle $\approx$ Erregungsrückbildung
\end{itemize}}

\Cave{Das EKG zeigt nur die \E{elektrische} Aktivität,
nicht aber die tatsächliche Muskelarbeit an! 

\EE{Es erlaubt
keine Aussage über die Auswurfleistung! }}

\Fazit{Grundsatz: \EE{Selbst ein pulsloser Patient kann ein
"`schönes"' \Abk{EKG} haben (allerdings nicht lange)!}}

\subsubsection{EKG -- Ableitungen}
\label{topic:ekg-ableitungen}

\Layout{0}{\TxBeschreibung}{
Grundsätzlich unterscheidet man zwischen den Extremitätenableitungen
und den Brustwandableitungen. 

\begin{itemize}
    \item Extremitätenableitungen\Z{: \EKG{I}, \EKG{II}, \EKG{III}, \EKG{aVL}, \EKG{aVR}, \EKG{aVF}}
    \item Brustwandableitungen\Z{: \VonBis{\EKG{V1}}{\EKG{V6}}}
\end{itemize}
}{
\begin{itemize}
    \item Extremitätenableitungen
    \item Brustwandableitungen
\end{itemize}
}{}{}{}

Es haben sich folgende Begriffe eingebürgert:

\begin{description}
    \item[Extremitätenableitung] "`\T{Rhythmusstreifen}"'
    zur Beurteilung des Herzrhythmus und zum Monitoring
    \item[Extremitätenableitung\,+\,Brustwandableitung]
    (\T{12-Kanal-EKG}"') -- Damit können Schäden der
    jeweiligen Herzregion zugeordnet werden. Wenn vorhanden,
    ist das 12-Kanal-\Abk{EKG} (\emph{"`12er"'}) für die
    Diagnostik das Mittel der Wahl. 
\end{description}

\Layout{0}{Extremitätenableitungen}
{
Die Extremitätenableitungen werden mittels vier Elektroden
abgeleitet. Je eine Elektrode wird an den Enden der Extremitäten angebracht. Zur
Vereinfachung können die Elektroden auch am Übergang vom Rumpf zu
den Extremitäten angebracht werden (z.\,B. statt dem rechten Arm an der
rechten Schulter). Wichtig ist, dass die 
Elektrodenpärchen immer auf \EE{gleicher Höhe} angebracht
werden!

Aus den Informationen der vier Elektroden erzeugt das
EKG-Gerät \EE{sechs Ableitungen} \Z{(\EKG{I}, \EKG{II}, \EKG{III}, \EKG{aVL}, \EKG{aVR} und \EKG{aVF}).}
}{
\begin{itemize}
    \item 4 Elektroden an den Extremitäten
    
     (bzw. am Übergang Rumpf -- Extremitäten).
    \item 6 Ableitungen.
\end{itemize}
}{}{}{}





\begin{figure}[H]
\begin{center}
\includegraphics[width=7cm]{./images/geraete/ECGcolor.png}
\Captionof{figure}{\AuthNotesPicLegalOk{}{PD}Extremitätenableitungen
(schwarze Kabel) und Brustwandableitungen (blaue Kabel).
\SourcePicture{WmCo/Madhero88, PD}\AuthNotes{GABS: Bessere Beschriftung wäre
wünschenswert, Dringlichkeit: niedrig}}
\end{center}
\end{figure}


\begin{table}[H]
\FormatTable
\Captionof{table}{Farbe und Positionen der EKG-Elektroden
für die Extremitätenableitung.}

\begin{tabular}{lll}
\TabHeading{Elektrode} 
& 
\TabHeading{Position} 
&
\TabHeading{Alternative Position}
\\\midrule
\colorbox{red}{\sffamily\bfseries\textcolor{white}{Rot}} & Rechte Hand & Rechte
Schulter
\\
\colorbox{yellow}{\sffamily\bfseries Gelb} & Linke Hand & Linke
Schulter
\\
\colorbox{green}{\sffamily\bfseries\textcolor{white}{Grün}} & Linker Fuß  &
Linke Leiste
\\
\colorbox{black}{\sffamily\bfseries\textcolor{white}{Schwarz}} & Rechter Fuß &
Rechte Leiste \\\bottomrule
\end{tabular}
\end{table}

\Layout{0}{Brustwandableitungen}
{
Die Brustwandableitungen werden mittels sechs Elektroden
abgeleitet und erlauben eine genauere Zuordnung von Störungen zu der jeweiligen
Region des Herzens. Aus den Informationen der sechs
Elektroden erzeugt das \Abk{EKG}-Gerät \EE{sechs
Ableitungen}: \Z{\VonBis{\EKG{V1}}{\EKG{V6}}}. 
}{
\begin{itemize}
    \item 6 Elektroden am Brustkorb.
    \item 6 Ableitungen.
\end{itemize}
}{}{}{}





\begin{table}[H]
\FormatTable
\Z{
\Captionof{table}{Position der EKG-Elektroden bei der Brustwandableitung.}

\begin{tabulary}{\linewidth}[H]{CL}
\TabSub{Elektrode}
 &
\TabSub{Position}
\\\midrule
\EKG{V1} & 4. Zwischenrippenraum rechts vom Sternum
\\
\EKG{V2} & 4. Zwischenrippenraum links vom Sternum
\\
\EKG{V3} & Mitte zwischen \EKG{V2} und \EKG{V4}
\\
\EKG{V4} & 5. Zwischenrippenraum links in der Verlängerung der
Schlüsselbein-Mittellinie
\\
\EKG{V5} & Mitte zwischen \EKG{V4} und \EKG{V6}
\\
\EKG{V6} & 6. Zwischenrippenraum in der mittleren Achsellinie
\\\bottomrule
\end{tabulary}
}
\end{table}


Die richtige Positionierung der \Abk{EKG}-Elektroden ist die Voraussetzung um
ein sinnvoll auswertbares \Abk{EKG} zu erstellen. Falsch angebrachte
Elektroden können ein \Abk{EKG} derart verfälschen, dass es zu
kompletten Fehldiagnosen kommen kann\ZFootnote{In \cite{Schnelle200607}
wurde untersucht, wie sich die Anbringungsstelle der Elektroden auf das
\Abk{EKG} auswirken kann. }.





\section{Körperliche Untersuchungen}

\subsection{Abtasten des Abdomens}

\Layout{0}{Abtasten des Abdomens}
{\label{topic:untersuchung-abdomen-abtasten}
Das Abdomen kann man abtasten um die Anspannung der
Bauchdecke zu beurteilen, Schmerzen zu
beurteilen und zuordnen, um Fremdkörper zu ertasten und
anderes. Man beginnt damit, den Patienten zu fragen, ob und wo er 
Schmerzen hat. Die Untersuchung
beginnt man in jenem Quadranten, welcher am weitesten von den
Schmerzen entfernt ist!

Man legt nun eine Hand flach auf den Bauch des Patienten und
die andere Hand flach auf die erste. Mit den Fingern drückt
man nun vorsichtig mehrere Zentimeter auf die Bauchdecke und
beurteilt, ob der Bauch weich oder verhärtet ist. Dabei
fragt man den Patienten ob er Schmerzen empfindet. Dies
wiederholt man in allen vier Quadranten.

Weiters kann man beurteilen ob der Patient Schmerzen durch
den Druck hat (Druckschmerz) oder ob er Schmerzen beim
Loslassen bekommt
(\T{Loslassschmerz}\index{Loslassschmerz}).
}
{
\begin{itemize}
    \item Fremdkörper? Wunden?
    \item Schmerzen?
    \item Einteilung in 4 Quadranten
    \item Abtasten der 4 Quadranten, beginnen dort wo kein
    Schmerz
    \item harter oder weicher Bauch?
    \item Druck- oder Loslassschmerz?
\end{itemize}
}{}{}{}

\clearpage % TODO temp clearpage
\subsection{Neurocheck}
\label{topic:neurocheck}
\index{Neurocheck}
%\begin{multicols}{2}

\Definition{Einfache, überblickshafte ("`grobe"')
Untersuchung der wichtigsten neurologischen Funktionen.}

\Layout{0}{Neurocheck}{
\begin{enumerate}
  \item \EE{Bewusstseinslage}: Für eine grobe Einschätzung
  ist folgende Differenzierung ausreichend:
  	\begin{enumerate}
        \item bewusstseinsklar
        \item bewusstseinsgetrübt
        \item bewusstlos
      \end{enumerate}
      Notärzte und Notfallsanitäter verwenden zur
      Beurteilung des Bewusstseins eine detailliertere
      Skala, die so genannte Glasgow Coma Scale.
  \item \EE{Orientierung}:
  	\begin{enumerate}
        \item zur Person (Patient kennt seinen Namen)
        \item zum Ort (Patient weiß wo er sich befindet)
        \item zur Zeit (Patient kann das ungefähre Datum
        und die ungefähre Uhrzeit nennen)
    \end{enumerate}
  \item \EE{Pupillen} Begutachtung und Lichtreaktion
\begin{enumerate}
  \item Pupillen gleich weit?
  \item Blickstarre (\T{Herdblick})? \index{Herdblick}
  \item \EE{Lichtreaktion:} In jedes Auge 2\,x leuchten und
  jeweils das gleiche und das andere Auge beobachten
  
  Beurteilung von:
  \begin{enumerate}
      \item Orientierung zu Person, Ort, Zeit und Situation.
  \item Schnelligkeit der Reaktion: prompt, verzögert?
  \item Reagieren immer beide
  Pupillen?~\Z{\footnote{\Z{isokor}}}
    \end{enumerate}
\end{enumerate}
  \item Lichtscheu und Sehstörungen?
  \item \EE{Kraftprobe} grobe Prüfung an oberer und
  unterer Extremität, beurteilt wird die Kraft und die
  Seitengleichheit. 
    \begin{enumerate}
  \item Patient überkreuzt die Hände geben, bei kräftigen
  Patienten nur Zeige- und Mittelfinger.
  \item Patienten bitten kräftig zuzudrücken.
  \item Mit Handflächen von unten gegen die Zehen drücken
  und Patienten bitten \emph{"`auf das Gaspedal zu
  drücken"'}.
  \item Von oben auf die Zehen drücken und Patienten bitten
  die Zehen Richtung Kopf zu heben.
\end{enumerate}
  \item \EE{Nackensteifigkeit}\label{topic:nackensteifigkeit-untersuchung} (ein
  {Meningismuszeichen}\index{Meningismuszeichen!Nackensteifigkeit!Untersuchung}\index{Nackensteifigkeit!Untersuchung})
  Der Patient liegt \EE{flach} auf dem Rücken. Der
  Untersucher fasst den Kopf des Patienten und versucht das Kinn an die
  Brust zu führen. Es handelt sich um eine \EE{passive
  Bewegung}, der Patient darf dabei nicht mithelfen.
  Normalerweise ist diese Bewegung ohne Probleme möglich.
  
  Wenn eine Nackensteifigkeit vorliegt ist die Beugung nicht
  möglich und der Versuch schmerzhaft, der Patient bewegt
  evtl. den ganzen Rumpf mit. Die Bedeutung der
  Nackensteifigkeit wird unter \prettyref{topic:meningitis}
  besprochen. \EE{Bei jedem Patienten mit Fieber ist auf das
  Vorliegen einer Nackensteifigkeit zu untersuchen.}  
  \item \EE{Blutzuckermessung} Eine Blutzuckermessung soll
  grundsätzlich bei einem Neurocheck erfolgen.
\end{enumerate}
}{
\begin{itemize}
  \item Bewusstseinslage
  \item Orientierung
  \item Pupillen
  \item Lichtscheue, Sehstörungen
  \item Kraftprobe
  \item Nackensteifigkeit
  \item[+] Blutzuckermessung
\end{itemize}
}{}{}{}

%\end{multicols}

\Layout{57}{}
{}{}
{
\begin{tabularx}{\linewidth}{XXX}
\includegraphics[width=\linewidth]{./images/motal-aass/00800/neuro1.jpg}
&
\includegraphics[width=\linewidth]{./images/motal-aass/00800/neuro2.jpg}
&
\includegraphics[width=\linewidth]{./images/motal-aass/00800/neuro3b.jpg}
\\
Pupillenreaktion, Lichtscheue, Sehstörungen 
&
Kraftprobe an den Armen
& 
Kraftprobe an den Beinen
\\\addlinespace
\includegraphics[width=\linewidth]{./images/motal-aass/00800/neuro4.jpg}
&
\includegraphics[width=\linewidth]{./images/motal-aass/00800/neuro5.jpg}
&
\includegraphics[width=\linewidth]{./images/motal-aass/00800/neuro6.jpg}
\\
Nackensteifigkeit
&
Nackensteifigkeit
&
Blutzuckermessung
\\\addlinespace\bottomrule
\end{tabularx}
}{Bilderserie: Neurocheck}{}

% FEHLT: Bild für Orientierung (stattdessen jetzt 2x
% Nackensteifigkeit)

\clearpage % TODO temp clearpage
\subsection{Traumacheck}
\label{topic:traumacheck}
\index{Traumacheck}
\Definition{Suche nach Verletzungen oder Verletzungsanzeichen am (liegenden)
Patienten.}

% Man unterscheidet zwischen einem \EE{ausführlichen} (vollständigen) und einem \EE{schnellen}
% Traumacheck. Der schnelle Traumacheck wird umfasst nur die wichtigsten Punkte,
% und wird in der Praxis oft vorgezogen, um rasch zu erkennen, ob der Patient
% vital bedroht ist.
%  
% \Layout{0}{\Ca{Schneller} Traumacheck}
% {%
% \begin{itemize}
%     \item Inspektion und Abtasten von
% \begin{enumerate}
%     \item Kopf,
%     \item Hals,
%     \item Thorax inkl. Atemprobe,
%     \item Bauch,
%     \item Becken und 
%     \item Oberschenkel
% \end{enumerate} 
% \item Kontrolle von Mund, Nase und Ohren auf Blutungen (mit einer
%     Lampe)
% \item Kann der Patient \E{Hände} und \E{Füße} spüren und bewegen?
% % KOCH: komplettes DMS würd zu lange dauern, Schuhe ausziehen, etc. 
% % GABS: Rekap-Zeit ist eh im normalen Ersteinschätzungsblock dabei.
% % MOTM: POI 35: "`ja klingt verünftig."'
% \end{itemize}
% }
% {
% \begin{itemize}
%     \item Rasche Ersteinschätzung
%     \item Rasches Erkennen von schweren, lebensbedrohlichen Verletzungen
% \end{itemize}
% {\SymbolText}
% }{}{}{}
% 
% 
% \Layout{0}{\E{Ausführlicher} Traumacheck}
% {%
% \begin{itemize}
%   \item Patient orientiert? Unfallhergang erinnerlich? Schmerzen?
%   \item Kontrolle von Mund, Nase und Ohren auf Blutungen (mit einer
%   Lampe).
%   \item Inspektion und Abtasten mit festem Griff, \EE{bei Schmerzen die vom Patienten erwähnte
%   Stelle als letztes untersuchen}: von Kopf bis Fuß
%   \begin{itemize}
%     \item Schädel: Schädeldach, Hinterkopf,
%     Stirn, Wangenknochen, Nasenbein, Kiefer.
%     \item Hals.
%     \item Obere Extremität: Schultern, Arme, Hände.
%     \item Brustkorb: vorne und seitlich (Prellmarken? Atemmechanik?, gegen leichten
%     Druck einatmen lassen).
%     \item Bauch in 4 Quadranten (Prellmarken?
%     Abwehrspannung/harte Bauchdecke?),
%     \item Becken (stabil?), Schambein.
%     \item Untere Extremität: Beine (sichtbare Verletzungen?), Füße:
%     Schuhe und Socken ausziehen)
%   \end{itemize}
%   \item Durchblutung, Motorik, Sensibilität (\T{DMS})
%   \begin{itemize}
%     \item Finger und Zehen: Kontrolle von \E{Motorik} (bewegen lassen), \E{Sensibilität} ({\quotedblbase}Spüren
%     Sie das?{\textquotedblleft}) und \E{Durchblutung}:
%     \item \T{Nagelbettprobe} (\EE{Rekapillarisierungszeit}): leichter Druck auf
%     Nagelende bis Nagelbett weiß ist, loslassen. Zeit bis Nagelbett wieder
%     rosarot ist sollte nicht länger als 1-2 Sekunden sein.
%     Seitenvergleich! Eine verlängerte Rekapzeit deutet auf eine
%     Durchblutungsstörung hin. (einseitig oder beidseitig? lokal oder systemisch?)
%   \end{itemize}
%   \item \T{Tupfertest}\index{Tupfertest}\label{topic:tupfertest}: etwas Blut aus
%   Ohren oder Nase mit einem Tupfer auffangen. Sollte Liquor im Blut enthalten sein bildet sich um den roten Blutfleck
%   ein heller, bernsteinfarbener Hof.
% \end{itemize}
% Bei Hinweis auf eine Schädelverletzung ist ein Neurocheck auf jeden Fall
% erforderlich!
% }{
% \begin{itemize}
%   \item Bewusstsein, Orientierung
%   \item Abtasten
%   \begin{itemize}
%     \item Schmerzen
%     \item Blutungen
%     \item Fehlstellungen
%     \item DMS
%     \item Atemmechanik
%     \item harte Bauchdecke
%   \end{itemize} 
%   \item Kontrolle von Mund, Nase, Ohren
%   \item Tupfertest
% \end{itemize}
% }{}{}{}


Teile des Traumachecks werden bereits während des
Einschätzungsblockes
(\CO{[PAM]}, \FullPrettyRef{topic:einschaetzungsblock})
durchgeführt.

\Layout{0}{\E{Vollständiger} Traumacheck}
{%
\begin{itemize}
  \item Patient orientiert? Unfallhergang erinnerlich? Schmerzen?
  \item Kontrolle von Mund, Nase und Ohren auf Blutungen (mit einer
  Lampe).
  \item Inspektion und Abtasten mit festem Griff, \EE{bei Schmerzen die vom Patienten erwähnte
  Stelle als letztes untersuchen}: von Kopf bis Fuß
  \begin{itemize}
    \item Schädel: Schädeldach, Hinterkopf,
    Stirn, Wangenknochen, Nasenbein, Kiefer.
    \item Hals.
    \item Obere Extremität: Schultern, Arme, Hände.
    \item Brustkorb: vorne und seitlich (Prellmarken? Atemmechanik?, gegen leichten
    Druck einatmen lassen).
    \item Bauch in 4 Quadranten (Prellmarken?
    Abwehrspannung/harte Bauchdecke?),
    \item Becken (stabil?), Schambein.
    \item Untere Extremität: Beine (sichtbare Verletzungen?), Füße:
    Schuhe und Socken ausziehen)
  \end{itemize}
  \item Durchblutung, Motorik, Sensibilität (\T{DMS})
  \begin{itemize}
    \item Finger und Zehen: Kontrolle von \E{Motorik}
    (bewegen lassen), \E{Sensibilität} (\Speech{"`Spüren Sie das?"'})
    und \E{Durchblutung}:
    \item \T{Nagelbettprobe} (\EE{Rekapillarisierungszeit}): leichter Druck auf
    Nagelende bis Nagelbett weiß ist, loslassen. Zeit bis Nagelbett wieder
    rosarot ist sollte nicht länger als \VonBis{1}{2} Sekunden sein.
    Seitenvergleich! Eine verlängerte Rekapzeit deutet auf eine
    Durchblutungsstörung hin. (einseitig oder beidseitig? lokal oder systemisch?)
  \end{itemize}
  \item \T{Tupfertest}\index{Tupfertest}\label{topic:tupfertest}: etwas Blut aus
  Ohren oder Nase mit einem Tupfer auffangen. Sollte Liquor im Blut enthalten sein bildet sich um den roten Blutfleck
  ein heller, bernsteinfarbener Hof.
\end{itemize}
Bei Hinweis auf eine Schädelverletzung ist ein
\E{Neurocheck} (\FullPrettyRef{topic:neurocheck}) auf jeden
Fall erforderlich! }{
\begin{itemize}
  \item Bewusstsein, Orientierung
  \item Abtasten
  \begin{itemize}
    \item Schmerzen
    \item Blutungen
    \item Fehlstellungen
    \item DMS
    \item Atemmechanik
    \item harte Bauchdecke
  \end{itemize} 
  \item Kontrolle von Mund, Nase, Ohren
  \item Tupfertest
\end{itemize}
}{}{}{}

\Layout{57}{}
{}{}
{
\begin{tabularx}{\linewidth}{XXX}
\includegraphics[width=\linewidth]{./images/motal-aass/00800/traumacheck1.jpg}
&
\includegraphics[width=\linewidth]{./images/motal-aass/00800/traumacheck2.jpg}
&
\includegraphics[width=\linewidth]{./images/motal-aass/00800/traumacheck3b.jpg}
\\
Kleidung entfernen 
&
Schuhe ausziehen (Fixierung, Stiefelgriff) 
& 
Untersuchung des Kopfes 
\\\addlinespace
\includegraphics[width=\linewidth]{./images/motal-aass/00800/traumacheck4.jpg}
&
\includegraphics[width=\linewidth]{./images/motal-aass/00800/traumacheck5b.jpg}
&
\includegraphics[width=\linewidth]{./images/motal-aass/00800/traumacheck6b.jpg}
\\
In alle Körperöffnungen im Kopf leuchten, Pupillenkontrolle.
&
Neben Blut kann auch Liquor austreten (Tupfertest).
&
Untersuchung des Thorax (inkl. Atemmechanik)
\\\addlinespace
\includegraphics[width=\linewidth]{./images/motal-aass/00800/traumacheck7b.jpg}
&
\includegraphics[width=\linewidth]{./images/motal-aass/00800/traumacheck8b.jpg}
&
\includegraphics[width=\linewidth]{./images/motal-aass/00800/traumacheck9b.jpg}
\\
Untersuchung des Abdomen und des Beckens
& 
Untersuchung der Extremitäten (DMS, Nagelbettprobe) 
&
Nach dem Traumacheck den entkleideten Patient vor
Wärmeverlust schützen. 
\\\addlinespace\bottomrule
\end{tabularx}
}
{Bilderserie: Traumacheck. Achte auf den Wärmeerhalt! In
kühler Umgebung soll ein ausführlicher Traumacheck erst
im (geheizten) Fahrzeug durchgeführt werden.}
{}






\section*{\protect\dsliterary{} Weiterführende Literatur}
\begin{WidePar}
\begin{multicols}{2}
\begin{labeling}{mmm}
  \item[\cite{DahmerAnamneseBefund10}] \emph{\bibentry{DahmerAnamneseBefund10}}
  \item[\cite{LehmeyerAnamUnter1}] \emph{\bibentry{LehmeyerAnamUnter1}}
  \item[\cite{ChecklisteAnamUnt1}] \emph{\bibentry{ChecklisteAnamUnt1}}
  \item[\cite{BatesUntersuchung1}] \emph{\bibentry{BatesUntersuchung1}}
\end{labeling}
\end{multicols}
\end{WidePar}






